\section{问题重述与分析}
微网通过外网购电、光伏利用与储能时移满足小区负载。购电计划不足会触发五倍电价紧急补购，过量合同又存在未消纳成本，因此经济性取决于储能调度与信息不确定性的联合处理。依据题目及附件\cite{problem}，四问分别研究固定负载与光伏预报下的确定性计划、未知净负荷下的日前计划、光伏预报更新下的日内调整，以及波动电价下前两类策略的变化。

本文以相同储能物理模型为基础：问题一建立确定性线性规划；问题二利用历史净负荷场景近似未知分布，以日前费用和期望紧急费用联合优化；问题三围绕最新预报叠加历史预测误差，在已执行状态下滚动调整；问题四仅扩展价格维度，以保证比较可解释。评价同时考察实际费用、紧急补购规模及尾部风险，避免仅用求解器目标值判断策略效果。

\section{模型假设与符号说明}
\begin{enumerate}
\item 将每条十分钟功率记录视作该段代表功率，忽略段内变化；光伏小时预报用线性插值在十分钟中点取值。附件标签与内部时段映射差异单独说明。
\item 采用每日首末储电量均为6000 kWh的周期运行边界。该条件强于题目给出的1月1日初值，排除了跨日能量转移；1月仅用于历史窗口初始化，未模拟其逐日储能轨迹。
\item 外网可按题设倍率补足实际缺口，不另设紧急购电容量上限；供能富余不向外网出售。忽略未给定的电池退化费用和设备启停成本。
\item 前31天负载与同日光伏或预测误差的联合变化可近似代表目标日不确定性。极端天气和突发负荷可能超出有限样本覆盖范围。
\item 第四问将附件4目标日完整价格曲线视为给定条件。结论对应已知价格曲线下的调度，未构建未来价格预测模型。
\end{enumerate}
\begin{table}[H]\centering\caption{主要符号}\small
\begin{tabular}{clc}\toprule 符号&含义&单位\\\midrule
$d,t,s$&日期、十分钟时段、历史场景&--\\
$p_{d,t}$&目标日期和交付时段的电价&元/kWh\\
$L,V,N$&负载、光伏及净负荷电量，$N=L-V$&kWh\\
$G,C,D$&合同购电、交流侧充电输入及放电输出&kWh\\
$E$&储电量，百分比SOC为$E/12000$&kWh\\
$U,W$&紧急购电及供能富余&kWh\\
$A^r,B^r$&发布时刻$r$的增购和削减量&kWh\\
$\pi_s$&场景概率，取$1/31$&--\\
$F_d,K_d^r$&实际日费用、一次调整净费用&元\\\bottomrule
\end{tabular}\end{table}

\section{数据预处理与共同约束}
\subsection{数据范围与信息边界}
附件1提供144条电价、负载及光伏预测记录；附件2提供365天实际负载与光伏，附件3每日提供四次未来24小时光伏预报，附件4提供365乘144的电价矩阵。按附件5规定，问题二至四输出2025年2月1日至12月31日334天，1月初始化31天历史窗口。检查数据维数、日期连续性、非负性及有限值，并保持原始行列顺序，不作小时聚合。

对每段取$\Delta t=1/6$小时，$L=P^{\rm L}\Delta t$、$V=P^{\rm PV}\Delta t$。实际数据只在决策冻结后用于结算；第三问在更新时仅额外使用已发布预报和最近已完成记录的光伏观测。加载全年文件本身不等于使用未来信息，关键在于求解器收到的截断场景与当前可用预报。

\subsection{时间标签与统计口径}
内部第$t$段表示从$(t-1)/6$到$t/6$小时，滚动更新在完成36、72个记录后进行。附件记录标记从0:10到0:00+1，模板指定区间按标签逐位置映射；例如模板10:00--10:10查询第60条，而该条在内部结束时刻约定下对应9:50--10:00。四小时统计始终按连续24条记录求和。本研究保留官方模板的既有位置映射，不平移数值，指定区间表与内部物理区间的边界解释仍存在十分钟差异。以下经济性比较均在相同记录位置下进行，不能将这些表理解为统一时标后重新优化的结果。

\subsection{储能物理约束}
充、放电分别按交流侧计量，效率均为0.9，故往返效率为0.81：
\begin{align}
E_t&=E_{t-1}+0.9C_t-D_t/0.9,\\
1200\leq E_t&\leq10800,\quad 0\leq C_t,D_t\leq5000\Delta t,\\
E_0&=E_{144}=6000.\label{eq:boundary}
\end{align}
外网购电、供能富余及紧急购电均非负。日内滚动时初始储电量由已执行轨迹继承，不能重置为6000。各问目标均加入$\varepsilon\sum(C+D)$，$\varepsilon=10^{-8}$，用于减少冗余吞吐；报告的实际费用剔除该项。小惩罚不等价于严格的充放电互斥约束，需事后核验。

\subsection{评价指标}
主要报告合同费用、紧急购电量与费用、实际总费用、发生紧急购电的天数以及日费用尾部指标。紧急电量占比是补购电量与实际外网购电总量之比，不能直接视作负载满足概率。因缺口由外网补足，本模型中的紧急事件表示计划偏差风险，不表示最终停电。供能富余可能含已支付的合同电量，不能全部解释为弃光。
